\section{Structure catégorielle du TPK : états enrichis, morphismes de
transport et classification opératoire}
\label{sec:arbol-operatorio-tpk-rev11}
\index{TPK!classification opératoire}
\index{sélecteur de phase!sous-opérateur du TPK}

Le TPK est le conteneur catégorique qui transporte sur APP les états
orientés de TRIT. Sa définition précède l'inventaire des réalisations :
on fixe d'abord les objets, les morphismes, la composition et les classes opératoires ; ce n'est
qu'ensuite que l'on introduit les opérateurs qui réalisent chaque classe.

\HMTClaim{HMT-R-TPK-ENRICHED-STATE}
\begin{definition}[TPK comme catégorie sur la base observable avec famille de fibres]
\label{def:estado-enriquecido-tpk-rev11}
Soit \(\mathsf P_{\rm obs}\) la catégorie des routes observables sur le graphe
de Cayley d'APP. Pour chaque configuration \(q\), on fixe la fibre
\[
 \mathcal F_q=
 \mathsf O_q\times\mathsf H_q\times\mathsf C_q\times
 \mathsf S_q\times\mathsf B_q\times\mathsf\Lambda_q,
\]
dont les facteurs conservent, respectivement, l'orientation et le feuillet ; le résidu et le
quotient ; la retenue ; la signature ; le cylindre et la frontière ; et le registre de transport. Un
objet du TPK est une paire \(x=(q,f)\), avec \(f\in\mathcal F_q\). Un morphisme
\(\Gamma:x\to y\) est une route \(\gamma:q\to q'\) de
\(\mathsf P_{\rm obs}\) accompagnée du transport de fibre qui préserve les
relations d'extension déclarées entre \(f\) et \(f'\). L'identité est la
route vide avec transport identité et la composition est la concaténation des
routes avec composition des transports de fibre. La projection
\[
 p:\mathcal X_{\rm TPK}^{\rm enr}\longrightarrow\mathsf P_{\rm obs}
\]
envoie chaque état sur sa configuration observable et chaque morphisme sur la route qu'il
transporte. Cette définition fixe une catégorie sur la base et une famille de
fibres ; elle ne présuppose pas l'existence de relèvements cartésiens pour toute
flèche de \(\mathsf P_{\rm obs}\).
\end{definition}

\begin{definition}[Classification ontologico--opératoire du TPK]
\label{def:categorias-operatorias-tpk-rev11}
La structure interne du TPK s'organise en huit classes différenciées par le type de
domaine, de codomaine et de fonction :
\begin{enumerate}
 \item \textbf{Support et état.} La base est \(\mathsf P_{\rm obs}\) et le
 porteur est \(\mathcal X_{\rm TPK}^{\rm enr}\to\mathsf P_{\rm obs}\). Cette
 classe contient des objets ; les sept suivantes contiennent des applications entre objets
 ou des lecteurs de ces objets.
 \item \textbf{Sélection interne.} La fonction de phase
 \[
 M_{\rm ph}:\{1,\ldots,9\}\longrightarrow\{+1,-1,0\}
 \]
 produit la valeur scalaire que l'opérateur
 \(\operatorname{Sel}_t\) reçoit pour activer l'évaluation additive,
 multiplicative ou le passage du seuil. C'est \(\operatorname{Sel}_t\), et non
 \(M_{\rm ph}\) isolément, qui agit sur la coordonnée de feuillet ; aucun
 des deux ne modifie à lui seul la position dans APP.
 \item \textbf{Transport.} Pour chaque direction cardinale \(d\),
 \(T_d:X_9\to X_9\), \(p\mapsto p+\delta(d)\), transporte la position.
 L'évolution \(F_{\rm obs}:\mathcal Q_{\rm obs}\to\mathcal Q_{\rm obs}\)
 compose sélection, translation, orientation et avancement de phase.
 \item \textbf{Lecteurs et quotients.} \(\rho_9\) publie la classe nonadique ;
 \((r_3,c_3)\) conserve le résidu tritique et le quotient ; et
 \((q_{1000},r_{1000})\) sépare le bloc décimal et la retenue. Ce sont des applications de
 lecture, non des transitions de l'état.
 \item \textbf{Mémoire et frontière.} Les transports de
 \(\mathsf H_q,\mathsf C_q,\mathsf S_q,\mathsf B_q\) et de
 \(\mathsf\Lambda_q\) mettent à jour le quotient, la retenue, la signature, le cylindre, la frontière
 et le registre. Leur projection visible peut coïncider alors que les états
 enrichis diffèrent.
 \item \textbf{Périodicité et retour.} L'extension
 \[
 0\longrightarrow C_{12}\longrightarrow C_{108}
 \longrightarrow C_9\longrightarrow0
 \]
 type le secteur, le calendrier et la phase visible. Le retour nonadique restaure la
 phase ; la périodicité de cent huit itérations restaure le calendrier ;
 aucun n'efface
 par définition la mémoire de frontière. Le lacet
 \(\mathcal K_{27}=F_{\rm obs}^{27}\) est un morphisme composé, non une autre
 horloge.
 \item \textbf{Prolongation.} Les relations
 \(\operatorname{Ext}_r\subseteq\mathcal F_q\times\mathcal F_{q'}\)
 sélectionnent les continuations compatibles. Les cylindres et leur système inverse
 publient la survie finie et, sous les hypothèses de compatibilité
 déclarées, une section globale.
 \item \textbf{Sorties.} Les foncteurs d'évaluation contractent le même état
 enrichi vers la sortie archimédienne, la sortie d'incidence
 exceptionnelle ou les classificateurs de l'atlas. Ces sorties sont des codomaines de
 lecture ; elles ne constituent pas de nouveaux conteneurs aux côtés du TPK.
\end{enumerate}
\end{definition}

Chaque opérateur qui apparaît ci-dessous est déterminé par le sextuplet
\((D,C,f,I,L,M)\) : domaine \(D\), codomaine \(C\), règle \(f\), invariants
préservés \(I\), information publiée ou perdue \(L\) et mémoire conservée
\(M\). L'égalité des cardinaux, des périodes ou des noms historiques n'identifie pas
deux opérateurs ayant des sextuplets distincts. Une fois cette
classification fixée, le diagramme suivant résume cinq regroupements fonctionnels
de réalisations qui raffinent les huit classes opératoires du TPK :
\[
\begin{tikzpicture}[
  node distance=7mm and 9mm,
  every node/.style={font=\small,align=center},
  root/.style={draw=hmtblue,very thick,rounded corners,fill=hmtlightblue,
    inner sep=5pt},
  op/.style={draw=hmtblue,rounded corners,fill=white,inner sep=4pt},
  sub/.style={draw=hmtgray,rounded corners,fill=hmtpaper,inner sep=3pt},
  edge/.style={-{Latex[length=2mm]},semithick,hmtblue}
]
\node[root] (tpk) {TPK\\$\mathcal X_{\rm TPK}^{\rm enr}$};
\node[op,below left=of tpk] (sel) {selección\\de lectura};
\node[op,below=of tpk] (trans) {transporte\\y retorno};
\node[op,below right=of tpk] (lift) {elevación\\y memoria};
\node[sub,below=of sel] (mov) {$\operatorname{Sel}_t[\,M_{\rm ph}(g(t))\,]$};
\node[sub,below=of trans] (card) {$T_d,\ \mathcal K_{27},\ F_{\rm obs}$};
\node[sub,below=of lift] (fib) {$\mathsf H,\mathsf C,\mathsf S,
  \mathsf B,\mathsf\Lambda$};
\node[op,below=17mm of trans] (clock) {calendario y registro\\
  $C_{12}\hookrightarrow C_{108}\twoheadrightarrow C_9$, $\Theta_{108}$};
\node[op,below=of clock] (ext) {prolongación, cilindros, frontera\\
  y filtración de supervivencia};
\draw[edge] (tpk) -- (sel);
\draw[edge] (tpk) -- (trans);
\draw[edge] (tpk) -- (lift);
\draw[edge] (sel) -- (mov);
\draw[edge] (trans) -- (card);
\draw[edge] (lift) -- (fib);
\draw[edge] (mov) |- (clock);
\draw[edge] (card) -- (clock);
\draw[edge] (fib) |- (clock);
\draw[edge] (clock) -- (ext);
\end{tikzpicture}
\]
Le diagramme fixe une inclusion, non une succession de théories. En particulier,
le sélecteur interne, les translations cardinales, l'observation
stroboscopique et les lecteurs modulaires conservent les types fixés dans
\cref{def:categorias-operatorias-tpk-rev11}. Le TPK coordonne leurs compositions
et conserve leur provenance dans un même état enrichi.

\HMTClaim{HMT-R-TPK-OPERATOR-SEPARATION}
\begin{proposition}[Séparation de la sélection, du mouvement et de l'observation]
\label{prop:tpk-separacion-operadores-rev11}
La sélection \(\operatorname{Sel}_t[\,M_{\rm ph}(g(t))\,]\), le transport
cardinal \(\operatorname{Lift}(T_d)\) et l'échantillonnage
\(\operatorname{Obs}_3\) sont des opérations non interchangeables : la première choisit
l'évaluation active ; la deuxième modifie la position et l'enroulement ; la
troisième publie une sous-suite temporelle. Leur composition au sein du TPK
détermine une route observable, tandis que la fibre \(\mathcal F_q\) conserve
les données nécessaires pour distinguer des prolongations ayant la même extrémité.
\end{proposition}

\begin{proof}
Les types sont différents : \(M_{\rm ph}\) prend ses valeurs dans l'ensemble ternaire
d'activation, \(\operatorname{Sel}_t\) met à jour la coordonnée de feuillet,
\(\operatorname{Lift}(T_d)\) transporte l'état sur le tore et
\(\operatorname{Obs}_3\) agit sur des suites. Une identification entre deux
de ces objets violerait le typage du domaine ou du codomaine. L'extension par fibres
incorpore en outre le résidu--quotient, la retenue, la frontière et le registre, de sorte
que la projection sur l'extrémité du chemin est une application d'oubli explicite.
\end{proof}

\begin{statusbox}[title={Nomenclature canonique}]
La dénomination publique du conteneur est \emph{TPK}. Lorsqu'il sera nécessaire de
nommer la donnée transportée, on emploiera \emph{état enrichi du TPK} ou
\(\mathcal X_{\rm TPK}^{\rm enr}\). La fonction \(M_{\rm ph}\) désigne
exclusivement la valeur interne de phase que reçoit
\(\operatorname{Sel}_t\).
\end{statusbox}
